%% =========================================================================
%%  Введение
%% =========================================================================

\chapter*{Введение}
\addcontentsline{toc}{chapter}{Введение}

\section*{Актуальность темы}
\addcontentsline{toc}{section}{Актуальность темы}

Квазиньютоновские методы занимают центральное место в вычислительной
оптимизации с момента работы Broyden (1965)~\cite{broyden1965} и
последующих классических работ Dennis и More~\cite{dennis1977,dennis1996}.
Они дают сверхлинейную скорость сходимости без вычисления якобиана
(или гессиана), что особенно важно для задач, в которых явное
дифференцирование дорого или недоступно: крупномасштабных систем
уравнений, задач обучения нейронных сетей, равновесных и экономических
моделей. На геометрически сложных ландшафтах (изогнутые долины, овраги,
неквадратичные сильно нелинейные функции) поведение классических
алгоритмов Broyden, SR1 и BFGS, однако, оказывается чувствительным
к выбору начального
приближения~\cite{more1981,dennis1996,conn1991,shannophua1978}, хотя
теория гарантирует сверхлинейность в локальном
режиме~\cite{broyden1973,dennis1974}.

Стандартный инструмент анализа сходимости квазиньютоновских
методов~--- оценка ограниченной деградации
(bounded~deterioration~\cite{broyden1973,dennis1996}), описывающая
рост нормы матрицы ошибки $E_k = B_k - J(x^*)$ через проектор вдоль
текущего шага. У классических обновлений ранга~1 (Бройден, SR1)
и~ранга~2 (BFGS) этот проектор имеет ранг~1: из~$p$ прошлых секущих
условий $B_{k+1}s_j = y_j$, $j = k{-}p,\ldots,k{-}1$, на~следующую
итерацию переходит лишь текущее. Естественный путь усиления
оценки~--- сохранение прошлых секущих условий, расширяющее проектор
ограниченной деградации с~ранга~1 до~ранга $p{+}1$, как для нелинейных
систем, так и для гладкой оптимизации. В настоящей диссертации такие
обновления построены и теоретически обоснованы: SP-Broyden для систем
$F(x) = 0$ и секантно-стабилизированное обновление SS-$\mathcal{U}$
для гладкой минимизации.

\section*{Объект и предмет исследования}
\addcontentsline{toc}{section}{Объект и предмет исследования}

\textbf{Объект исследования} --- численные методы решения двух
взаимосвязанных классов задач: нелинейных систем уравнений
$F(x) = 0$ и безусловной гладкой оптимизации $\min_{x\in\RR^n} f(x)$.

\textbf{Предмет исследования} --- глобальные квазиньютоновские
обновления приближённого якобиана $B_k$, сохраняющие на каждой
итерации $p$ прошлых секущих условий $B_{k+1}s_j = y_j$,
$j=k{-}p,\ldots,k$, и порождающие расширенный (ранг $p{+}1$)
проектор bounded deterioration; в частности, ранг-1 обновление
SP-Broyden для нелинейных систем и его симметричный аналог~SS-U
для оптимизации.

\section*{Цель и задачи работы}
\addcontentsline{toc}{section}{Цель и задачи работы}

\textbf{Цель работы:} построить и теоретически обосновать
квазиньютоновские обновления приближённого якобиана с памятью
$p$ прошлых секущих пар $(s_j, y_j)$ и применить полученные
алгоритмы к задачам нелинейных систем и гладкой оптимизации.

\textbf{Задачи:}
\begin{enumerate}[leftmargin=*]
  \item Построить алгоритм SP-Broyden для нелинейных систем:
        выбор направления обновления, адаптивное управление глубиной
        окна по числу обусловленности грамиана, процедура обновления.
  \item Доказать для SP-Broyden теорему об ограниченной деградации
        с расширенным проектором: точное разложение
        $\|E_{k+1}\|_F^2 = \|E_k\|_F^2 - \|E_k\Pi_k\|_F^2 + \Delta_k$
        с проектором $\Pi_k$ ранга $p{+}1$ (у классического Бройдена
        --- ранга~1) и неравенство
        $\|E_k\Pi_k\|_F^2 \ge \|E_k s_k\|^2/\|s_k\|^2$.
  \item Перенести идею сохранения прошлых секущих условий
        на симметричные методы (SR1, PSB): отталкиваясь от
        классической несовместности симметрии и сохранения
        нескольких секущих условий~\cite{schnabel1983}, построить
        секантно-стабилизированные обновления~SS-U с коррекцией в
        ортогональном дополнении $s_k^{\perp}$, исследовать
        сверхлинейность и устойчивую сходимость на модельных задачах.
  \item Подтвердить теоретические результаты численными экспериментами
        на стандартных тестовых задачах.
\end{enumerate}

\section*{Научная новизна}
\addcontentsline{toc}{section}{Научная новизна}

\begin{enumerate}[leftmargin=*]
  \item Предложено и исследовано семейство обновлений~SP-Broyden~---
        ранг-1 минимальное обновление, сохраняющее $p$ прошлых секущих
        условий $B_{k+1}s_j = y_j$, $j = k{-}p, \ldots, k{-}1$.
        Для него получено \emph{точное разложение} bounded deterioration
        $\|E_{k+1}\|_F^2 = \|E_k\|_F^2 - \|E_k\Pi_k\|_F^2 + \Delta_k$
        с проектором $\Pi_k$ ранга $p{+}1$ (у классического Бройдена~---
        ранга~1) и локальной оценкой
        $\Delta_k = O(\rho_k^2)$, где
        $\rho_k = \max_{k-p\le j\le k}\|x_j-x^*\|$.
        Доказано неравенство
        $\|E_k\Pi_k\|_F^2 \ge \|E_k s_k\|^2/\|s_k\|^2$~---
        отрицательный член точного разложения для SP-Broyden
        не меньше соответствующего члена классического Бройдена.
  \item Для симметричного случая, отталкиваясь от классической
        несовместности симметрии и одновременного сохранения нескольких
        секущих условий~\cite{schnabel1983}, предложена конструкция~SS-U
        --- ранг-1 коррекция в~$s_k^{\perp}$, не нарушающая ни текущего
        секущего условия, ни симметрии и стабилизирующая невязку прошлых
        секущих в смысле минимизации в~$s_k^{\perp}$. Установлена
        сверхлинейность~SS-U по Денниса--Море и численно продемонстрирована
        устойчивая сходимость на модельных задачах.
\end{enumerate}

\section*{Теоретическая и практическая значимость}
\addcontentsline{toc}{section}{Теоретическая и практическая значимость}

\textbf{Теоретическая значимость.} Получено точное разложение
$\|E_{k+1}\|_F^2 = \|E_k\|_F^2 - \|E_k\Pi_k\|_F^2 + \Delta_k$
для семейства~SP-Broyden~--- равенство, обобщающее классическую
оценку сверху bounded deterioration с одномерного проектора
$s_k s_k^\top/\|s_k\|^2$ на проектор ранга $p{+}1$ на
подпространство прошлых шагов; отрицательный член для SP-Broyden
не меньше соответствующего члена классического Бройдена.
В симметричной квазиньютоновской теории, отталкиваясь от
классической несовместности симметрии и одновременного сохранения
нескольких секущих условий~\cite{schnabel1983}, предложена
конструкция~SS-U как ранг-1 коррекция в~$s_k^\perp$, разрешающая
это препятствие; для неё установлена Q-сверхлинейность по критерию
Денниса--Море.

\textbf{Практическая значимость.} Limited-memory вариант
L-SP-Broyden хранит $O(nm)$ векторов вместо плотной матрицы
$O(n^2)$ и применим к нелинейным системам высокой размерности,
где плотные квазиньютоновские формы исчерпывают память. По числу
итераций до сходимости L-SP-Broyden не уступает плотному
SP-Broyden при $m\ge p{+}1$ и не деградирует с ростом размерности
на стандартных тестовых задачах; адаптивное управление глубиной
окна по числу обусловленности грамиана обеспечивает устойчивость
к плохо обусловленным шагам без ручной настройки параметров.
По сравнению с классическим Бройденом расширенное окно секущих
сокращает медианное число итераций до сходимости и сужает
разброс по случайным стартам. В оптимизационной части семейство
SS-U восстанавливает устойчивую сходимость симметричных
обновлений~PSB и~SR1 на задачах, где базовое обновление
структурно деградирует, не ухудшая сходимость на задачах, где
деградации нет.

\section*{Положения, выносимые на защиту}
\addcontentsline{toc}{section}{Положения, выносимые на защиту}

\begin{enumerate}[leftmargin=*]
  \item Характеризация ранг-1 обновлений якобиана, сохраняющих
        $p+1$ прошлых секущих условий, и явная формула SP-Broyden
        как такого обновления.
  \item Теорема минимальности обновления SP-Broyden: единственное
        решение задачи минимизации нормы Фробениуса при
        $(p{+}1)$~секущих условиях (см.~теорему~\ref{thm:minimality}).
  \item Теорема об ограниченной деградации с расширенным проектором
        для SP-Broyden (см.~теорему~\ref{thm:bd}): точное разложение
        ошибки $\|E_{k+1}\|_F^2 = \|E_k\|_F^2 - \|E_k \Pi_k\|_F^2 + \Delta_k$
        с проектором $\Pi_k$ ранга $p{+}1$ (у классического Бройдена~---
        ранга~1) и неравенство
        $\|E_k\Pi_k\|_F^2 \ge \|E_k s_k\|^2/\|s_k\|^2$.
  \item Конструкция секантно-стабилизированного обновления~SS-U
        для симметричного случая: ранг-1 коррекция в~$s_k^{\perp}$,
        разрешающая классическую несовместность симметрии и
        одновременного сохранения нескольких секущих
        условий~\cite{schnabel1983}.
  \item Q-сверхлинейная сходимость SS-$\mathcal{U}$ по критерию
        Денниса--Море (см.~теорему~\ref{thm:ss_dm}).
\end{enumerate}

\section*{Методология и методы исследования}
\addcontentsline{toc}{section}{Методология и методы исследования}

В работе используются методы вычислительной линейной алгебры
(разложения QR и SVD, теория проекторов, оценки чисел обусловленности
грамиана) и классической теории квазиньютоновских методов
(свойства обновлений ранга~1 и~2, критерий сверхлинейности
Dennis--More~\cite{dennis1977}).
Численные эксперименты проведены в~Python с использованием~NumPy;
задачи взяты из стандартного тестового набора~\cite{more1981}.

\section*{Степень достоверности и апробация результатов}
\addcontentsline{toc}{section}{Степень достоверности и апробация результатов}

Все теоретические результаты снабжены полными доказательствами.
Утверждения о скорости сходимости подтверждены численными экспериментами
на стандартных тестовых задачах с единым критерием остановки.
Исходный код экспериментов и скрипты построения всех представленных
рисунков открыты и доступны в публичном
репозитории~\texttt{SP-Broyden}.\footnote{\url{https://github.com/mojaevr/secant-conditions}}

\section*{Структура и объём работы}
\addcontentsline{toc}{section}{Структура и объём работы}

Работа состоит из введения, трёх глав, заключения и списка литературы.

\begin{itemize}[leftmargin=*]
  \item В \textbf{главе~1} (<<Приближение якобиана>>) вводится язык
        квазиньютоновских обновлений и формулируется общее семейство
        мультисекантных обновлений. Анализируется выбор матрицы
        направлений $V_m$ и базовой матрицы $\widehat{B}_k$.
  \item В \textbf{главе~2} (<<Алгоритм SP-Broyden>>) строится обновление SP-Broyden,
        доказываются теоремы минимальности и ограниченной деградации,
        описывается адаптивный выбор глубины $p$ и приводятся
        численные эксперименты для нелинейных систем.
  \item В \textbf{главе~3} идея сохранения прошлых секущих условий
        переносится на симметричные методы. Строится SS-U,
        доказывается сверхлинейная сходимость, исследуется
        устойчивая сходимость на модельных задачах.
\end{itemize}
